% !TeX program = xelatex
% ============================================================================
%  第 9 课　正方形里的小群：子群与二面体群　—— 学生版讲义
%  与 02-课件/第09课-正方形的小群/第09课-正方形的小群.tex 同步
% ============================================================================

\xhlesson{9}{正方形里的小群：子群与二面体群}{}

\begin{mubiao}
  \begin{itemize}
    \item 能找出 \(D_4\) 里面能自成一体的小群（子群）。
    \item 能说出旋转与翻折之间那条关系：\(fr=r^3f\)（先翻后转 \(=\) 先倒着转再翻）。
    \item 知道 \(D_n\) 叫二面体群，正 \(n\) 边形有 \(2n\) 个对称动作。
  \end{itemize}
\end{mubiao}

\section*{一、找 \(D_4\) 的子群}

\noindent 办法照旧：圈出一小堆动作，把复合填成一张表。全在里面，就是子群。

\begin{dongshou}[先检查两个]
  \begin{enumerate}
    \item \(H=\{e,\ r^2\}\)：
      \(r^2\cdot r^2=\) \underline{\hspace{2em}}，跑出去了吗？\underline{\hspace{3em}}
      \quad 所以它\underline{\hspace{2em}}（是／不是）子群。
    \item \(H=\{e,\ r\}\)：
      \(r\cdot r=\) \underline{\hspace{2em}}，这个结果在 \(H\) 里吗？\underline{\hspace{3em}}
      \quad 所以它\underline{\hspace{2em}}（是／不是）子群。
  \end{enumerate}
  对一下答案：\(\{e,r^2\}\) \textbf{是}子群；\(\{e,r\}\) \textbf{不是}——\(r\cdot r=r^2\) 跑出去了。
\end{dongshou}

\begin{dongshou}[再填一张 \(4\) 阶的表：\(H=\{e,\ r^2,\ f,\ r^2f\}\)]
  横竖都按 \(e,\ r^2,\ f,\ r^2f\) 排。格子里的动作还是「\textbf{先横排的，再做竖排的}」。
  四行都填，看看有没有哪个格子跑出这四个动作。

  \begin{center}
  \begin{tabular}{@{}c|cccc@{}}
    \toprule
     & \(e\) & \(r^2\) & \(f\) & \(r^2f\) \\
    \midrule
    \rule{0pt}{3ex}\(e\) & & & & \\
    \midrule
    \rule{0pt}{3ex}\(r^2\) & & & & \\
    \midrule
    \rule{0pt}{3ex}\(f\) & & & & \\
    \midrule
    \rule{0pt}{3ex}\(r^2f\) & & & & \\
    \bottomrule
  \end{tabular}
  \end{center}
  跑出去了吗？\underline{\hspace{3em}}\quad 所以它\underline{\hspace{2em}}（是／不是）子群，大小是 \underline{\hspace{1.2em}}。
\end{dongshou}

\begin{example}[\(D_4\) 里全部的子群]
  \begin{center}
  \begin{tabular}{@{}clc@{}}
    \toprule
    大小 & 子群 & 个数 \\
    \midrule
    \(1\) & \(\{e\}\) & \(1\) \\
    \(2\) & \(\{e,r^2\}\)；\(\{e,f\}\)；\(\{e,rf\}\)；\(\{e,r^2f\}\)；\(\{e,r^3f\}\) & \(5\) \\
    \(4\) & \(\{e,r,r^2,r^3\}\)；\(\{e,r^2,f,r^2f\}\)；\(\{e,r^2,rf,r^3f\}\) & \(3\) \\
    \(8\) & \(D_4\)（全部八个） & \(1\) \\
    \bottomrule
  \end{tabular}
  \end{center}
  一共 \(10\) 个。
\end{example}

\begin{sikao}
  看最大子群和最小子群的大小：\(1,2,4,8\)。它们和 \(D_4\) 的大小 \(8\) 有什么关系？\\[4pt]
  （和上节课 \(S_3\) 里的发现比一比。）
  \liukong{3}
\end{sikao}

\section*{二、旋转与翻折的关系}

\begin{dongshou}[左右各走一条路]
  先把两条路写成两行式（提示：\(r\) 是转 \(90^\circ\)，\(r^3\) 是转 \(270^\circ\)，
  \(f\) 是沿对角线 \(1\text{—}3\) 翻）：
  \[
    fr=\begin{pmatrix}1&2&3&4\\ \hspace{1.2em}&\hspace{1.2em}&\hspace{1.2em}&\hspace{1.2em}\end{pmatrix},
    \qquad
    r^3f=\begin{pmatrix}1&2&3&4\\ \hspace{1.2em}&\hspace{1.2em}&\hspace{1.2em}&\hspace{1.2em}\end{pmatrix}
  \]
  两次的终点一样吗？\quad 一样 \(\square\)\qquad 不一样 \(\square\)
\end{dongshou}

\begin{example}
  \[
    fr=r^3f .
  \]
  白话：「先翻后转，等于先\textbf{倒着}转（转 \(270^\circ\)）、再翻。」
  \(r^3\) 就是 \(r\) 的\textbf{逆元}——把 \(r\) 转的 \(90^\circ\) 倒回去。
\end{example}

\begin{definition}[二面体群]
  正 \(n\) 边形（\(n\ge 3\)）的全部对称动作，按复合构成一个群，叫做\textbf{二面体群}，记作 \(D_n\)，
  一共有 \(2n\) 个动作：\(n\) 个转动加 \(n\) 个翻折。
\end{definition}

\noindent 提醒一句：三角形的 \(D_3\)，就是我们前面一直写的 \(S_3\)——同一批 \(6\) 个动作，
数学家只是给了它两个名字。

\begin{sikao}
  \(D_3\) 有多少个动作？\(D_5\)（正五边形）呢？\(D_{12}\) 呢？
  \liukong{2}
\end{sikao}

\begin{xiaojie}
  猜一猜：只用 \(r\) 和 \(f\) 反复复合，能不能造出全部 \(8\) 个动作？（下节课揭晓）
  \liukong{3}
\end{xiaojie}

\noindent\textbf{下节课}：生成元、循环群与循环群的子群——为什么两个就够。